\documentclass[10pt,a4paper]{article}

\usepackage[utf8]{inputenc}
\usepackage[T1]{fontenc}
\usepackage[ngerman]{babel}
\usepackage[a4paper,top=1.0cm,bottom=1.2cm,left=0.8cm,right=0.8cm]{geometry}
\usepackage{amsmath,amssymb}
\usepackage[table]{xcolor}
\usepackage{multicol}
\usepackage{fancyhdr}
\usepackage{array}

% ---------- Farben ----------
\definecolor{head}{rgb}{0.00,0.45,0.55}   % Hauptueberschrift (kräftiges Cyan)
\definecolor{sub}{rgb}{0.00,0.58,0.68}    % Unterueberschrift (helleres Cyan)
\definecolor{grau}{rgb}{0.45,0.45,0.45}   % Hinweise

% ---------- Layout / Abstaende ----------
\setlength{\parindent}{0pt}
\setlength{\parskip}{0.4pt}
\setlength{\columnsep}{10pt}
\setlength{\abovedisplayskip}{1.5pt}
\setlength{\belowdisplayskip}{1.5pt}
\setlength{\abovedisplayshortskip}{0.5pt}
\setlength{\belowdisplayshortskip}{0.5pt}
\renewcommand{\baselinestretch}{0.90}

% ---------- Ueberschriften ----------
\newcommand{\Sec}[1]{\par\vspace{1.8pt}{\color{head}\normalsize\bfseries #1}\nopagebreak\par\vspace{-1pt}{\color{head}\rule{\linewidth}{0.9pt}}\par\nopagebreak\vspace{1pt}}
\newcommand{\Sub}[1]{\par\vspace{1pt}{\color{sub}\bfseries #1}\nopagebreak\par\vspace{0.2pt}}
\newcommand{\hint}[1]{{\color{grau}\itshape #1}}

% ---------- Symbole ----------
\newcommand{\jc}{{\color{sub}\boldmath$\checkmark$}}   % ja  (Haken)
\newcommand{\nc}{{\color{grau}\boldmath$\times$}}      % nein (Kreuz)

% ---------- kompakte Aufzaehlung ----------
\newenvironment{tl}{\begin{list}{\textbullet}{%
  \setlength{\leftmargin}{1.1em}\setlength{\labelsep}{0.45em}%
  \setlength{\itemsep}{0pt}\setlength{\parsep}{0pt}\setlength{\topsep}{1pt}%
  \setlength{\partopsep}{0pt}}}{\end{list}}

% ---------- Kopf-/Fusszeile ----------
\pagestyle{fancy}
\fancyhf{}
\renewcommand{\headrulewidth}{0pt}
\renewcommand{\footrulewidth}{0.4pt}
\fancyfoot[L]{\footnotesize\color{grau}Optimierungsverfahren — Open-Book-Formelsammlung}
\fancyfoot[R]{\footnotesize\color{grau}Seite \thepage}

\begin{document}
\footnotesize

\begin{center}
{\color{head}\Large\bfseries Optimierungsverfahren — Formelsammlung}\\[1pt]
{\color{grau}\footnotesize Grundbegriffe \,·\, Modellierung \,·\, Lineare Optimierung \& Transport \,·\, Hesse-Matrix \& Karush-Kuhn-Tucker \,·\, Gradientenfreie Verfahren \,·\, Metaheuristiken \,·\, Mehrzieloptimierung \,·\, Methoden-Tabelle}
\end{center}
\vspace{2pt}

\begin{multicols}{3}

% ============================================================
\Sec{Grundbegriffe der Optimierung}

\begin{tl}
\item \textbf{unimodal:} nur ein Optimum \quad \textbf{multimodal:} mehrere Optima (lokale und globale)
\item \textbf{lokaler Algorithmus:} löst nur unimodale Probleme
\item \textbf{globaler Algorithmus:} löst auch multimodale Probleme
\end{tl}
\[ \text{lokal: } \arg\min_{x\in N(x_0)} f \qquad \text{global: } \arg\min_{x\in\Omega} f \]
\hint{$N(x_0)$ Umgebung, $\Omega$ ganzer Suchraum. Bei konvexen Problemen: lokal $=$ global.}

% ============================================================
\Sec{Modellierung}

\[ \min_{\vec x} f(\vec x)\ \ \text{unter}\ \ h_k(\vec x)=0,\ g_j(\vec x)\le 0,\ x_i^l\le x_i\le x_i^u \]
\begin{tl}
\item Entscheidungsvariablen (mit Einheit!)
\item Zielfunktion $f$: was wird minimiert oder maximiert?
\item Nebenbedingungen: Gleichungen $h$, Ungleichungen $g$, Grenzen
\end{tl}
\textbf{Linear} genau dann, wenn $f$ \emph{und} alle Nebenbedingungen linear in $\vec x$ sind \hint{(kein $\sqrt{\cdot},\,x^2,\,\sin,\,x_ix_j$)}.

% ============================================================
\Sec{Lineare Optimierung: Standardform}

Solver verlangen eine einheitliche Form:
\begin{tl}
\item nur Minimierung
\item nur Gleichheits-Nebenbedingungen
\item nur nicht-negative Variablen ($x_i\ge 0$)
\item rechte Seite stets positiv ($b_j\ge 0$)
\end{tl}
\[ \min_{\vec x} \vec c^{\,\top}\vec x \quad\text{unter}\quad A\vec x=\vec b,\ \ \vec x\ge 0 \]
Ungleichung $\to$ Gleichung über Slackvariablen.

% ============================================================
\Sec{Simplex-Algorithmus}

\begin{tl}
\item Start an einem Eckpunkt des zulässigen Bereichs
\item entlang der Kante mit der größten Verbesserung zum nächsten Eckpunkt gehen
\item beenden, wenn keine Verbesserung mehr möglich
\end{tl}
\hint{Im Kern wiederholtes Lösen linearer Gleichungssysteme (Gauß); Optimum liegt in einem Eckpunkt.}

% ============================================================
\Sec{Ganzzahlige lineare Optimierung}

Variablen müssen ganzzahlig sein ($\vec x\in\mathbb{Z}^n$).
\begin{tl}
\item \textbf{naiv:} kontinuierlich lösen, dann auf nächste zulässige ganze Zahl runden \hint{(nicht immer optimal)}
\item \textbf{Schnittebenenverfahren:} (1) kontinuierlich lösen; (2) ganzzahlig? dann fertig; (3) sonst einen Schnitt hinzufügen, der die nicht-ganzzahlige Lösung ausschließt, alle ganzzahligen aber erlaubt; (4) wiederholen
\end{tl}

% ============================================================
\Sec{Grafische Lösung (2 Variablen)}

\Sub{Linear}
\begin{tl}
\item Nebenbedingungen als Geraden $\Rightarrow$ zulässiger Bereich (Schnitt der Halbebenen)
\item Isolinien von $f$ verschieben, Optimum meist in einer \textbf{Ecke}
\end{tl}

\Sub{Nichtlinear}
Zulässig sind die Punkte, die \emph{alle} Nebenbedingungen erfüllen (bei Gleichungs-Nebenbedingungen: \textbf{Schnittpunkte} der Kurven \emph{innerhalb} der Grenzen). Davon den mit bestem Wert von $f$ wählen; $\vec x$ und $f(\vec x)$ ablesen.

\Sub{Slackvariable $s_j\ge 0$}
\[ g_j(\vec x)+s_j=b_j \]
\textbf{aktive} Nebenbedingung (Gerade berührt) $\Rightarrow s_j=0$; \textbf{inaktive} $\Rightarrow s_j>0$.

% ============================================================
\Sec{Klassisches Transportproblem}

$m$ Werke (Angebot $a_i$) $\to$ $n$ Läden (Nachfrage $b_j$), Kosten $c_{ij}$, Mengen $x_{ij}\ge 0$.
\begin{tl}
\item \textbf{zulässig:} Zeilensumme $=a_i$, Spaltensumme $=b_j$
\item \textbf{klassisch/balanciert:} $\sum a_i=\sum b_j$
\item Kosten $=\sum_{i,j} c_{ij}\,x_{ij}$
\end{tl}

\Sub{Startlösung}
\begin{tl}
\item \textbf{Nord-West-Ecken-Regel:} nordwestlichste leere Zelle mit Minimum aus Angebot und Nachfrage füllen; Angebot/Nachfrage verringern; wiederholen \hint{(ignoriert Kosten)}
\item \textbf{Matrixminimum-Methode:} immer die Zelle mit den geringsten Kosten zuerst füllen \hint{(meist günstiger)}
\item \textbf{Vogel-Methode:} pro Zeile/Spalte Opportunitätskosten (zweitkleinste $-$ kleinste Kosten); höchste wählen, dort günstigste Zelle füllen
\end{tl}

\Sub{Optimale Lösung}
Verbesserung der Startlösung über \textbf{Stepping-Stone} oder \textbf{MODI}.

% ============================================================
\Sec{Stationäre Punkte \& Hesse}

\Sub{Notwendig}
\[ \nabla f(\vec x)=\vec 0\ \Rightarrow\ \text{stationäre Punkte} \]

\Sub{Hinreichend \hint{($H=\nabla^2 f$ im Punkt)}}
\begin{tl}
\item $H$ positiv definit $\Rightarrow$ Minimum
\item $H$ negativ definit $\Rightarrow$ Maximum
\item $H$ indefinit $\Rightarrow$ Sattelpunkt
\item semidefinit $\Rightarrow$ unentscheidbar
\end{tl}
Bei einer $2\times 2$-Matrix: positiv definit $\Leftrightarrow H_{11}>0$ \emph{und} $\det H>0$; negativ definit $\Leftrightarrow H_{11}<0$ \emph{und} $\det H>0$; indefinit $\Leftrightarrow \det H<0$.

% ============================================================
\Sec{Gradientenverfahren}

\[ \vec x_{k+1}=\vec x_k-\eta\,\nabla f(\vec x_k),\quad \eta>0 \]
\hint{$\eta$ zu groß $\to$ Divergenz, zu klein $\to$ langsam. Findet nur \textbf{lokales} Min, braucht Gradient.}

% ============================================================
\Sec{Lagrange \& Karush-Kuhn-Tucker}

Minimiere $f$ unter $h(\vec x)=0$ (Multiplikator $a$) und $g(\vec x)\le 0$ (Multiplikator $b$):
\[ L(\vec x,a,b)=f(\vec x)+a\,h(\vec x)+b\,g(\vec x) \]
\textbf{Bedingungen:}
\begin{tl}
\item Stationarität: $\nabla_{\vec x}L=\vec 0$
\item Zulässigkeit: $h=0,\ g\le 0$
\item Komplementärschlupf: $b\cdot g(\vec x)=0$
\item Vorzeichen: $b\ge 0$ \hint{(Minimierung mit $g\le 0$)}
\end{tl}

% ============================================================
\Sec{Quadratische Optimierung}

Quadratische Zielfunktion, lineare Nebenbedingungen:
\[ \min_{\vec x}\ \tfrac12\,\vec x^{\,\top} Q\,\vec x + \vec c^{\,\top}\vec x \quad\text{unter}\quad A\vec x\le\vec b,\ A_{eq}\vec x=\vec b_{eq} \]
\begin{tl}
\item $Q$ symmetrisch ($Q=Q^{\top}$)
\item $Q$ positiv definit $\Rightarrow$ konvex, eindeutige Lösung
\end{tl}
\hint{Verfahren: Innere-Punkte, konjugierter Gradient, Simplex-Varianten, erweiterte Lagrange-Methode.}

% ============================================================
\Sec{Sequentielle quadratische Optimierung}

Für nichtlineare Probleme mit Nebenbedingungen: in jeder Iteration eine \textbf{quadratische Näherung} lösen.
\[ \min_{\vec d}\ f(\vec x_i)+\nabla f(\vec x_i)^{\top}\vec d+\tfrac12\,\vec d^{\,\top} H\vec d \]
\hint{$H$ = Hesse-Matrix der Lagrange-Funktion (oft genähert, BFGS).}
\begin{tl}
\item Lösung des Teilproblems liefert Suchrichtung $\vec d$
\item Schrittweite $c>0$ per Liniensuche
\item $\vec x_{i+1}=\vec x_i+c\,\vec d$, bis Konvergenz
\end{tl}
\hint{lokal, gradientenbasiert, deterministisch, behandelt Nebenbedingungen.}

% ============================================================
\Sec{Nelder-Mead (Simplex)}

Gradientenfreies Verfahren. Simplex aus $n{+}1$ Punkten, sortiert von $x_0$ (bestes) bis $x_n$ (schlechtestes).
\[ x_c=\tfrac1n\!\!\sum_{i\ne n}\! x_i,\qquad x_r=x_c+\alpha(x_c-x_n) \]
\begin{tl}
\item $f(x_r)<f(x_0)$ \hint{(neuer Bester)} $\Rightarrow$ \textbf{Expansion}
\item $f(x_0)\le f(x_r)<f(x_{n-1})$ $\Rightarrow$ $x_r$ übernehmen
\item $f(x_r)\ge f(x_{n-1})$ $\Rightarrow$ \textbf{Kontraktion}
\item sonst $\Rightarrow$ \textbf{Schrumpfen} (zu $x_0$)
\end{tl}
\hint{lokal; für multimodale Probleme mit wiederholten Neustarts global machen.}

% ============================================================
\Sec{DIRECT \hint{(DIviding RECTangles)}}

Global, gradientenfrei, deterministisch; nur Grenzen (Box). Je Iteration:
\begin{tl}
\item Menge der \textbf{potentiell optimalen} Rechtecke bestimmen \hint{(auf der konvexen Hülle aus Größe gegen Güte)}
\item jedes davon \textbf{abtasten und dritteln} (größte Dimension teilen)
\item wiederholen bis zur maximalen Iterationszahl
\end{tl}
\hint{Größe $d_i$ = Abstand Mittelpunkt$\to$Ecke, Güte $=f$ im Mittelpunkt.}

% ============================================================
\Sec{Evolutionäre Algorithmen}

\textbf{Population} von Lösungskandidaten; pro Generation: Eltern auswählen $\to$ rekombinieren $\to$ mutieren $\to$ auswerten $\to$ Überlebende auswählen.
\begin{tl}
\item \textbf{Mutation:} kleine zufällige Änderung, etwa $\vec x+\vec\epsilon$ mit $\vec\epsilon\sim N(0,\sigma)$; bei Bitstrings ein Bit kippen
\item \textbf{Rekombination:} z.\,B. Mittelwert $(\vec x_1+\vec x_2)/2$ oder Crossover
\item \textbf{Fitness} = Zielfunktion $f(\vec x)$
\end{tl}
\hint{stochastisch, gradientenfrei, lokal (kleine Mutationen) und global (Population). Hyperparameter: Populationsgröße, Mutationsrate und -schrittweite.}

% ============================================================
\Sec{(1+1)-Evolutionsstrategie \hint{(1/5-Erfolgsregel)}}

Der Elternpunkt wird mutiert; die bessere Lösung überlebt. \textbf{Erfolg} (bei Minimierung): $f(x_{\text{next}})<f(x)$.\\
Erfolgsrate $p$ über das Fenster (Speicherlänge), Zielrate $1/5$:
\begin{tl}
\item $p>1/5 \Rightarrow \sigma\leftarrow\sigma\cdot\alpha$ \hint{(vergrößern)}
\item $p<1/5 \Rightarrow \sigma\leftarrow\sigma/\alpha$ \hint{(verkleinern)}
\item $p=1/5 \Rightarrow \sigma$ unverändert
\end{tl}
\hint{Rate einschließlich aktuellem Schritt: Erfolge geteilt durch Fensterlänge.}

% ============================================================
\Sec{CMA-ES}

\textbf{Covariance Matrix Adaptation Evolution Strategy:} Evolutionsstrategie, die eine fortlaufend angepasste \textbf{Kovarianzmatrix} verwendet, um die ideale Verteilung der Individuen im Suchraum anzunähern.
\hint{Stand der Technik für kontinuierliche Optimierung; global bei ausnutzbarer globaler Struktur, sonst Neustarts oder große Population.}

% ============================================================
\Sec{Simulated Annealing}

Inspiriert vom Abkühlen von Metall: hohe Temperatur $\to$ viele zufällige Änderungen, langsames Abkühlen $\to$ immer kleinere Änderungen.
\[ \text{Verschlechterung } \Delta f \text{ akzeptiert mit } \exp\!\Big(-\tfrac{\Delta f}{T}\Big),\quad T\downarrow \]
\hint{erlaubt das Verlassen lokaler Optima.}

% ============================================================
\Sec{Particle Swarm Optimization}

Inspiriert von Schwärmen (Vögel, Insekten). Jedes Teilchen ist ein Lösungskandidat; gute Lösungen beeinflussen, wohin sich schlechtere Lösungen bewegen \hint{(Informationsaustausch im Schwarm)}.

% ============================================================
\Sec{Ant Colony Optimization}

Schwarm-Verfahren: Ameisen hinterlassen \textbf{Pheromone} auf ihrem Weg und folgen stärkeren Spuren wahrscheinlicher; kurze Wege erhalten mehr Pheromone.
\hint{typisch für Routing-Probleme.}

% ============================================================
\Sec{Surrogatmodell-basierte Optimierung}

Für \textbf{teure} Zielfunktionen: ein günstiges \textbf{Surrogatmodell} (Modell aus maschinellem Lernen) ersetzt die teure Auswertung.
\begin{tl}
\item Startpunkte auswerten (Versuchsplanung)
\item Surrogatmodell anpassen
\item auf dem Modell vielversprechenden nächsten Punkt wählen
\item dort echte Zielfunktion auswerten, Modell aktualisieren, wiederholen
\end{tl}

% ============================================================
\Sec{Mehrzieloptimierung}

\Sub{Dominanz \hint{(alle Ziele minimieren)}}
$P$ \textbf{dominiert} $Q$: $P\le Q$ in \emph{allen} Zielen \emph{und} $<$ in mindestens einem.
\begin{tl}
\item \textbf{Pareto-optimal:} von keinem Punkt dominiert
\item \textbf{Pareto-Menge:} alle Pareto-optimalen Punkte; ihr Bild im Zielraum = \textbf{Pareto-Front}
\item \textbf{Non-Dominated-Sorting:} Rang 1 = alle nicht dominierten; diese entfernen, im Rest erneut die nicht dominierten = Rang 2, und so weiter
\end{tl}

\Sub{Crowding Distance}
Maß für die Diversität innerhalb einer Front. Pro Zielfunktion:
\begin{tl}
\item Punkte nach dieser Zielfunktion sortieren
\item Randpunkte (kleinster und größter Wert): Distanz $=\infty$
\item innere Punkte $i$:
\end{tl}
\[ \text{Distanz}\mathrel{+}=\frac{f_{i+1}-f_{i-1}}{f_{\max}-f_{\min}} \]
\hint{$f_{i-1},f_{i+1}$ = Werte des linken und rechten Nachbarn.} Über alle Zielfunktionen summieren; \textbf{große Distanz wird bevorzugt}.

\Sub{NSGA-II}
Evolutionärer Algorithmus für mehrere Ziele; Auswahl über Rang und Crowding Distance:
\begin{tl}
\item initiale Population
\item Eltern wählen (Rang, dann Crowding Distance)
\item rekombinieren, mutieren
\item Überlebende der nächsten Generation wählen (Rang, dann Crowding Distance)
\item wiederholen
\end{tl}

\Sub{Hypervolumen}
Vom Rand (Referenzpunkt) bis zur Front \glqq überdecktes\grqq{} Volumen; größer = bessere Front. Der \textbf{Hypervolumenbeitrag} eines Punktes ist sein Anteil daran \hint{(Randpunkte oft $\infty$)} — alternativer Tie-Break zur Crowding Distance.

\end{multicols}

% ============================================================
% Volle Breite: Methoden-Vergleichstabelle
% ============================================================
\vspace{2pt}
{\color{head}\normalsize\bfseries Eigenschaften der Optimierungsmethoden}\par
\vspace{-1pt}{\color{head}\rule{\linewidth}{0.9pt}}\par\vspace{2pt}

\footnotesize
\setlength{\tabcolsep}{3pt}
\renewcommand{\arraystretch}{1.2}
\begin{center}
\begin{tabular}{|>{\raggedright\arraybackslash}p{2.9cm}|c|c|c|c|c|c|c|c|c|}
\hline
\rowcolor{head!10}
\textbf{Methode} & \textbf{Global} & \shortstack{\textbf{Gradienten-}\\\textbf{basiert}} & \shortstack{\textbf{Neben-}\\\textbf{bedingungen}} & \textbf{Ziele} & \textbf{Linear} & \textbf{Ansatz} & \textbf{Genauigkeit} & \shortstack{\textbf{Determi-}\\\textbf{nistisch}} & \textbf{Variablen}\\
\hline
Stationäre Punkte / Hesse-Matrix & \jc & \jc & ohne & 1 & \nc & analytisch & exakt & \jc & $\mathbb{R}^n$\\
\hline
Gradientenverfahren & \nc & \jc & ohne & 1 & \nc & numerisch & ? & \jc & $\mathbb{R}^n$\\
\hline
Simplex (lineare Optimierung) & \nc & \jc & alle & 1 & \jc & numerisch & exakt & \jc & $\mathbb{R}^n,\mathbb{Z}^n$\\
\hline
Quadratische Optimierung & \nc & \jc & alle & 1 & teilweise & numerisch & ? & \jc & $\mathbb{R}^n$\\
\hline
Karush-Kuhn-Tucker & \jc & \jc & alle & 1 & \nc & analytisch & exakt & \jc & $\mathbb{R}^n$\\
\hline
Sequentielle quadratische Optimierung & \nc & \jc & alle & 1 & \nc & numerisch & ? & \jc & $\mathbb{R}^n$\\
\hline
Nelder-Mead & \nc & \nc & ohne & 1 & \nc & numerisch & näherungsweise & \jc & $\mathbb{R}^n$\\
\hline
DIRECT & \jc & \nc & Grenzen & 1 & \nc & numerisch & näherungsweise & \jc & $\mathbb{R}^n$\\
\hline
Evolutionäre Algorithmen & \jc\,(?) & \nc\,(?) & / & / & \nc & numerisch & näherungsweise & \nc & /\\
\hline
(1+1)-Evolutionsstrategie & \nc & \nc & ohne & 1 & \nc & numerisch & näherungsweise & \nc & $\mathbb{R}^n$\\
\hline
CMA-ES & \jc\,(?) & \nc & ohne & 1 & \nc & numerisch & näherungsweise & \nc & $\mathbb{R}^n$\\
\hline
Simulated Annealing & \jc & \nc & ohne & 1 & \nc & numerisch & näherungsweise & \nc & $\mathbb{R}^n,\mathbb{Z}^n$\\
\hline
Particle Swarm Optimization & \jc & \nc & ohne & 1 & \nc & numerisch & näherungsweise & \nc & $\mathbb{R}^n$\\
\hline
Ant Colony Optimization & \jc & \nc & ohne & 1 & \nc & numerisch & näherungsweise & \nc & $\mathbb{Z}^n$\\
\hline
Surrogatmodell-basierte Optimierung & \jc\,(?) & \nc\,(?) & / & / & \nc & numerisch & näherungsweise & / & /\\
\hline
NSGA-II & \jc & \nc & ohne & mehrere & \nc & numerisch & näherungsweise & \nc & $\mathbb{R}^n$\\
\hline
\end{tabular}
\end{center}
\vspace{-2pt}
{\color{grau}\footnotesize \jc\,=\,ja \quad \nc\,=\,nein \quad (?)\,=\,abhängig von Problem und Parametern \quad „teilweise``\,=\,linear nur in den Nebenbedingungen \quad „/``\,=\,nicht festgelegt}

\end{document}
